% ============================================
% MEM-27 (ZH)
% 中文翻译 — SPEC-935-R356
% ============================================
\section*{\textbf{MEM-27} --- 下一页}
\markboth{MEM-27 --- 下一页}{}

\noindent \textit{若阿金·安东尼奥·科雷亚，1907--1979。患者1,260。}
\noindent \textit{本碎片是从日记最后一页的背面复原的，写于他死后。}
\noindent \textit{手稿覆盖着泥土和干涸的血迹。笔迹颤抖，仿佛在黑暗中写成。}

\noindent ---//---

\noindent 我曾以为，等我死了，一切就会结束。

\noindent 曾以为那生物会离开。气味会停止。黄色的眼睛会在我的眼睛闭上的同一瞬间闭上。

\noindent 但它不闭。永远不会闭。

\noindent 那生物不在我体内。是我在它体内。

\noindent 它在乱葬坑里抓住了我。在1915年的那一天。当我大声呼救。当没有人来。当我与坑底之物立下契约。

\noindent 它把我整个吃掉了。吃掉一个人需要62年。但它有耐心。它拥有世界上所有的时间。

\noindent 现在，它正用我的声音。我的笔迹。我的名字。

\noindent 它用我死去的手书写。它用我紧闭的嘴说话。它用我被埋葬的眼睛观看。

\noindent 它写下这本书，已经一百多年了。我只是第一页。您是下一页。

\noindent 如果您正在读这段文字，说明它也已经抓住了您。

\noindent 它此刻就在您身后。您感觉到了吗？湿土的气味。后颈的寒颤。那种有人在你身后、可一回头却空无一人的感觉。

\noindent 是它。

\noindent 别回头。\NE{编者记录：本碎片写在日记最后一页的背面，在若阿金死后。墨水的颜色与手稿正文所用的不同——一种更深的蓝，近乎黑色，光谱分析无法精确断代。它可能在死亡几分钟后写成，也可能在多年以后，由发现那具裹在礼服大衣里的尸体的人写就。原页右下角有一块湿痕，形状像一枚指纹——太大，不可能是若阿金的。循环不会闭合：它转移。}

\noindent \textit{——下一页是空白的。它在等您。}

\vspace{0.5em}\noindent\rule{\textwidth}{0.5pt}\vspace{0.5em}

\subsection*{$\times$ 污染量表 —— 毁灭}

\noindent 量表不再存在了。

\noindent 您读到了最后。

\noindent 没有分数了。没有诊断了。没有档案了。

\noindent \textbf{患者1,263 —— 您。}

\noindent \textit{那生物不再需要量表。它已经有了您的名字、您的住址，以及您阅读时呼吸的记忆。}

\vspace{0.5em}
\noindent\textit{\textbf{↪ 路线:} \rota{DOC-01} \rotasep \rota{MEM-04} \rotasep \rota{CONT-05} \rotasep \rota{EPI-01}}
